% !TEX program = xelatex
\documentclass[12pt,a4paper,fontset=windows]{ctexart}
\usepackage{amsmath,amssymb}
\usepackage[a4paper,margin=2.2cm,top=2.6cm,bottom=2.4cm,headheight=15pt]{geometry}
\usepackage{booktabs}
\usepackage{array}
\usepackage{xcolor}
\usepackage{enumitem}
\usepackage{tikz}
\usetikzlibrary{trees}
\usepackage[most]{tcolorbox}
\usepackage{fancyhdr}
\usepackage[hidelinks,unicode]{hyperref}
\hypersetup{pdftitle={期末复习模拟卷 + 强化卷（附详细解答）}}

\pagestyle{fancy}
\fancyhf{}
\fancyhead[C]{\small 期末复习模拟卷 + 强化卷（附详细解答）}
\fancyfoot[C]{\small —~\thepage~—}
\renewcommand{\headrulewidth}{0.4pt}

\setlist[enumerate]{itemsep=1pt,topsep=3pt,leftmargin=2.2em}
\setlist[itemize]{itemsep=2pt,topsep=3pt,leftmargin=1.6em}

\newcommand{\FV}[1]{\ensuremath{\mathrm{FIRSTVT}(#1)}}
\newcommand{\LV}[1]{\ensuremath{\mathrm{LASTVT}(#1)}}
\newcommand{\FVname}{\ensuremath{\mathrm{FIRSTVT}}}
\newcommand{\LVname}{\ensuremath{\mathrm{LASTVT}}}
\newcommand{\FR}[1]{\ensuremath{\mathrm{FIRST}(#1)}}
\newcommand{\FO}[1]{\ensuremath{\mathrm{FOLLOW}(#1)}}
\newcommand{\FRname}{\ensuremath{\mathrm{FIRST}}}
\newcommand{\FOname}{\ensuremath{\mathrm{FOLLOW}}}
\newcommand{\ok}{\ensuremath{\checkmark}}
\newcommand{\bad}{\ensuremath{\times}}
\newcommand{\imp}{\ensuremath{\Rightarrow}}

\newtcolorbox{qbox}[1][题目]{%
  enhanced,breakable,colback=gray!8,colframe=gray!55!black,
  boxrule=0.6pt,arc=2pt,left=7pt,right=7pt,top=5pt,bottom=5pt,
  before skip=8pt,after skip=8pt,
  title={\bfseries #1},coltitle=black,colbacktitle=gray!25,fonttitle=\bfseries}
\newtcolorbox{warnbox}[1][本题易错]{%
  enhanced,breakable,colback=yellow!12,colframe=orange!75!black,
  boxrule=0.7pt,arc=2pt,left=7pt,right=7pt,top=5pt,bottom=5pt,
  before skip=8pt,after skip=8pt,
  title={\bfseries #1},coltitle=black,colbacktitle=yellow!35,fonttitle=\bfseries}
\newtcolorbox{tipbox}[1][要点提示]{%
  enhanced,breakable,colback=blue!5,colframe=blue!50!black,
  boxrule=0.7pt,arc=2pt,left=7pt,right=7pt,top=5pt,bottom=5pt,
  before skip=8pt,after skip=8pt,
  title={\bfseries #1},coltitle=black,colbacktitle=blue!15,fonttitle=\bfseries}
\newtcolorbox{infobox}{%
  enhanced,breakable,colback=gray!6,colframe=gray!45!black,
  boxrule=0.5pt,arc=1mm,left=7pt,right=7pt,top=5pt,bottom=5pt,
  before skip=6pt,after skip=8pt}

\begin{document}

\begin{center}
{\LARGE\bfseries 期末复习模拟卷 — 五类题型自测（附详细解答）}
\end{center}

\begin{infobox}
\textbf{用途}：期末冲刺自测卷。覆盖《期末复习专题》题型 1--5，共 11 题（$1+3+3+1+3$），每题附\textbf{具体解答}。建议限时 60--90 分钟先独立作答，再对照解答批改。\\
\textbf{配套}：《期末复习专题》·《分析设计题题库》·《应用题题库》·《02-文法和语言》·《05-语法分析-自底向上》\\
\textbf{判卷要点}：结论先行、依据齐全，对照《期末复习专题》\textbf{备考总纲}与专题附录A 检查清单。
\end{infobox}

\section{题型1：FIRSTVT/LASTVT 与算符优先分析（第 1 题）}

\begin{qbox}
{\large\bfseries 第 1 题}\quad 已知文法 $G[E]$：
\[ E \to E \vee T \mid T,\qquad T \to T \wedge F \mid F,\qquad F \to \neg F \mid (E) \mid a \]
（其中 $\vee$、$\wedge$、$\neg$、$($、$)$、$a$ 为终结符）。\\
(1) 求各非终结符的 \FVname、\LVname；(2) 构造算符优先关系表；(3) 判断 $G$ 是否算符优先文法，并说明理由。
\end{qbox}

\textbf{解 (0) 拓广文法}：$G'[Z]$：$Z \to \#E\#$（引入界符 $\#$，兼作栈底与输入结束标志）。

\textbf{解 (1) 集合计算}（规则 F1/F2/F3、L1/L2/L3 见《期末复习专题》题型1 知识点）

\FVname{}：
\begin{itemize}
\item \textbf{Z}：$Z \to \#E\#$ 右部以终结符 $\#$ 开头，用 F1 得 $\#$ \imp{} $\FV{Z} = \{\,\#\,\}$
\item \textbf{F}：$F \to \neg F$、$F \to (E)$、$F \to a$ 右部均以终结符开头，用 F1 得 $\neg$、$($、$a$ \imp{} $\FV{F} = \{\,\neg,\ (,\ a\,\}$
\item \textbf{T}：$T \to T \wedge F$ 是 F2 形（右部以非终结符开头、紧跟终结符 $\wedge$）得 $\wedge$；$T \to F$ 用 F3 得 $\FV{F} \subseteq \FV{T}$ \imp{} $\FV{T} = \{\,\wedge,\ \neg,\ (,\ a\,\}$
\item \textbf{E}：$E \to E \vee T$ 用 F2 得 $\vee$；$E \to T$ 用 F3 得 $\FV{T} \subseteq \FV{E}$ \imp{} $\FV{E} = \{\,\vee,\ \wedge,\ \neg,\ (,\ a\,\}$
\end{itemize}

\LVname{}：
\begin{itemize}
\item \textbf{Z}：$Z \to \#E\#$ 右部以终结符 $\#$ 结尾，用 L1 得 $\#$ \imp{} $\LV{Z} = \{\,\#\,\}$
\item \textbf{F}：$F \to a$、$F \to (E)$ 右部以终结符结尾，用 L1 得 $a$、$)$；\textbf{$F \to \neg F$ 右部以非终结符结尾、其前紧跟终结符 $\neg$，是 L2 形 \imp{} $\neg \in \LV{F}$}（本题最易漏的一条！）\imp{} $\LV{F} = \{\,\neg,\ ),\ a\,\}$
\item \textbf{T}：$T \to T \wedge F$ 是 L2 形得 $\wedge$；$T \to F$ 用 L3 得 $\LV{F} \subseteq \LV{T}$ \imp{} $\LV{T} = \{\,\wedge,\ \neg,\ ),\ a\,\}$
\item \textbf{E}：$E \to E \vee T$ 用 L2 得 $\vee$；$E \to T$ 用 L3 得 $\LV{T} \subseteq \LV{E}$ \imp{} $\LV{E} = \{\,\vee,\ \wedge,\ \neg,\ ),\ a\,\}$
\end{itemize}

\textbf{解 (2) 算符优先关系}（行 = 栈顶符号，列 = 输入头符号；$<$ 低于、$\doteq$ 同级、$>$ 高于）

逐条产生式扫描：
\begin{itemize}
\item \textbf{$F \to (E)$}：$(\,E$ 形 \imp{} $( \lessdot \FV{E}$；$E\,)$ 形 \imp{} $\LV{E} \gtrdot )$；$(\,E\,)$ 两终结符隔一个非终结符 \imp{} $( \doteq )$。
\item \textbf{$F \to \neg F$}：$\neg\,F$ 形 \imp{} $\neg \lessdot \FV{F}$。
\item \textbf{$T \to T \wedge F$}：$T\,\wedge$ 形 \imp{} $\LV{T} \gtrdot \wedge$；$\wedge\,F$ 形 \imp{} $\wedge \lessdot \FV{F}$。
\item \textbf{$E \to E \vee T$}：$E\,\vee$ 形 \imp{} $\LV{E} \gtrdot \vee$；$\vee\,T$ 形 \imp{} $\vee \lessdot \FV{T}$。
\item \textbf{$Z \to \#E\#$}：$\#\,E$ 形 \imp{} $\# \lessdot \FV{E}$；$E\,\#$ 形 \imp{} $\LV{E} \gtrdot \#$；$\#\,E\,\#$ 两终结符隔一个非终结符 \imp{} $\# \doteq \#$。
\end{itemize}

得关系表：

\begin{center}
\renewcommand{\arraystretch}{1.35}
\begin{tabular}{c|ccccccc}
\toprule
  & $\vee$ & $\wedge$ & $\neg$ & $($ & $)$ & $a$ & $\#$ \\
\midrule
$\vee$   & $>$ & $<$ & $<$ & $<$ & $>$ & $<$ & $>$ \\
$\wedge$ & $>$ & $>$ & $<$ & $<$ & $>$ & $<$ & $>$ \\
$\neg$   & $>$ & $>$ & $<$ & $<$ & $>$ & $<$ & $>$ \\
$($      & $<$ & $<$ & $<$ & $<$ & $\doteq$ & $<$ &   \\
$)$      & $>$ & $>$ &   &   & $>$ &   & $>$ \\
$a$      & $>$ & $>$ &   &   & $>$ &   & $>$ \\
$\#$     & $<$ & $<$ & $<$ & $<$ &   & $<$ & $\doteq$ \\
\bottomrule
\end{tabular}
\end{center}

\textbf{解 (3) 判定}：拓广后的 $G'$ 无 $\varepsilon$ 产生式、各产生式右部无相邻非终结符 \imp{} 是\textbf{算符文法}；表中每一对终结符（含 $\#$）之间至多只有一种优先关系（如 $($ 与 $)$ 之间只有 $\doteq$，$\vee$ 与 $($ 之间只有 $<$，$\#$ 与 $\#$ 之间只有 $\doteq$），无冲突 \imp{} \textbf{$G[E]$ 是算符优先文法}。

\begin{warnbox}
\begin{enumerate}
\item \textbf{$F \to \neg F$ 是 L2 形}（右部以非终结符结尾、其前紧跟终结符）\imp{} $\neg \in \LV{F}$。漏掉它，``$\neg$'' 行会少 3 个 $>$。$\neg$ 同时在 \FV{F} 与 \LV{F} 里是正常的。
\item \textbf{``$\neg$'' 行既有 $<$ 又有 $>$ 不是冲突}：冲突指\textbf{同一格}出现两种以上关系；不同格不同关系互不影响。
\item 判定结论两句话：先``是算符文法''，再``每对终结符至多一种关系、无冲突''，最后``故是算符优先文法''。
\item \textbf{表格要带 $\#$ 行、$\#$ 列}：由 $Z \to \#E\#$ 得 $\# \lessdot \FV{E}$、$\LV{E} \gtrdot \#$、$\# \doteq \#$。$\#$ 行只有 $<$ 与 $\doteq$，$\#$ 列只有 $>$ 与 $\doteq$，方向别写反。
\end{enumerate}
\end{warnbox}

\section{\texorpdfstring{题型2：语言 $\to$ 文法，并判断文法类型（第 2--4 题）}{题型2：语言→文法，并判断文法类型（第 2–4 题）}}

\begin{qbox}
{\large\bfseries 第 2 题}\quad 写一文法 $G$，使其语言为 $L=\{a^{2n} \mid n \ge 0\}$（偶数个 $a$ 组成的串，含空串 $\varepsilon$），并判断文法类型。
\end{qbox}

\textbf{解}：
\[
\begin{aligned}
S &\to aA \mid \varepsilon\\
A &\to aS
\end{aligned}
\]
\begin{itemize}
\item \textbf{正向}：$S \to \varepsilon$ 得 $n=0$；$S \to aA$、$A \to aS$ 两步增加两个 $a$，反复使用得任意偶数个 $a$。
\item \textbf{反向}：每轮必须``成对''使用 $S \to aA$ 与 $A \to aS$（$A$ 不可单独终止），故 $a$ 的总数恒为偶数。
\item \textbf{类型}：$S \to aA$、$A \to aS$ 都是右线性（$aB$ 形），$S \to \varepsilon$ 允许 \imp{} \textbf{3 型（正规文法）}；语言是正规语言。
\end{itemize}

\begin{warnbox}
若写 $S \to aaS \mid \varepsilon$，语言对，但``$aaS$''不是右线性形式（右线性要求右部为 $A \to aB \mid a$）\imp{} 该文法只是 \textbf{2 型}。判类型是判\textbf{文法}不是判语言！（见《期末复习专题》题型2 易错点 3。）
\end{warnbox}

\begin{qbox}
{\large\bfseries 第 3 题}\quad 写一文法 $G$，使其语言为 $L=\{a^n b^n c^m \mid n \ge 1,\ m \ge 0\}$，并判断文法类型。
\end{qbox}

\textbf{解}：
\[
\begin{aligned}
S &\to AB\\
A &\to aAb \mid ab\\
B &\to cB \mid \varepsilon
\end{aligned}
\]
\begin{itemize}
\item $A$ 生成 $a^n b^n$（$n \ge 1$），$B$ 生成 $c^m$（$m \ge 0$，含空），拼接即得语言。
\item \textbf{正向}：$A \Rightarrow a^n b^n$、$B \Rightarrow c^m$ \imp{} $a^n b^n c^m$；\textbf{反向}：任意合法串按结构唯一拆成``等数 $a$、$b$ 段 + $c$ 段''，分别由 $A$、$B$ 推出。
\item \textbf{边界}：$m = 0$ 时 $B \Rightarrow \varepsilon$，如 \texttt{ab}、\texttt{aabb} \ok；$n \ge 1$ 保证至少一对 $a$、$b$ \ok。
\item \textbf{类型}：$S \to AB$ 右部含两个非终结符 \imp{} 非 3 型；左部为单个非终结符 \imp{} \textbf{2 型}。语言非正规（$a^n b^n$ 需要计数配对，DFA 做不到）。
\end{itemize}

\begin{warnbox}
不要写成 $S \to aSb \mid ab \mid C$，$C \to cC \mid \varepsilon$：此时 $S \Rightarrow aSb \Rightarrow aCb$ 会推出 \texttt{acb}、\texttt{accb} 等 \textbf{$c$ 夹在 $a$、$b$ 之间}的非法串。$c$ 段必须独立成块（用 $S \to AB$ 拼接），或显式限制 $S \to aSb$ 时内层 $S$ 不得再用 $C$ 产生式。
\end{warnbox}

\begin{qbox}
{\large\bfseries 第 4 题}\quad 写一文法 $G$，使其语言为 $L=\{a^n b^m \mid n > m \ge 0\}$，并判断文法类型。
\end{qbox}

\textbf{解}：
\[
S \to aSb \mid aS \mid a
\]
\begin{itemize}
\item \textbf{结构读法}：$S \to a$ 是出口（1 个 $a$）；$aS$ 每次加 1 个 $a$；$aSb$ 每次加 1 个 $a$ 和 1 个 $b$。设 $aSb$ 用 $j$ 次、$aS$ 用 $k$ 次，则串为 $a^{1+k+j} b^{j}$，即 $n = m + 1 + k > m$（$m = j \ge 0$）。
\item \textbf{正向}：推得的串都是 $a$ 在前、$b$ 在后且 $n > m \ge 0$；\textbf{反向}：任给 $n > m \ge 0$，取 $j = m$、$k = n-m-1 \ge 0$，即可推出 $a^n b^m$。
\item \textbf{边界}：$m = 0$ 得 \texttt{a}、\texttt{aa}、\texttt{aaa}…（$n \ge 1$）\ok；最小``带 $b$''串是 \texttt{aab}（$n=2,\ m=1$）\ok；\texttt{ab}（$n=m=1$）不在语言中，推不出 \ok。
\item \textbf{类型}：$S \to aSb$ 右部非右线性 \imp{} \textbf{2 型}；语言非正规（需比较 $n$ 与 $m$ 的大小差，DFA 记不住）。
\end{itemize}

\section{\texorpdfstring{题型3：文法 $\to$ 语言，并判断文法/语言类型（第 5--7 题）}{题型3：文法→语言，并判断文法/语言类型（第 5–7 题）}}

\begin{qbox}
{\large\bfseries 第 5 题}\quad $G[S]$：$S \to aS \mid bA \mid \varepsilon$，$A \to bA \mid \varepsilon$。求 $L(G)$，并判断文法类型与语言类型。
\end{qbox}

\textbf{解}：$S \to aS$ 在串前加一个 $a$，$S \to bA$ 转入``加 $b$''阶段（至少 1 个 $b$），$S \to \varepsilon$、$A \to \varepsilon$ 是出口。归纳可得：
\[
L(G)=\{a^i b^j \mid i,\ j \ge 0\} = a^* b^*
\]
（即任意个 $a$ 后随任意个 $b$，含 $\varepsilon$，如 \texttt{$\varepsilon$}、\texttt{a}、\texttt{b}、\texttt{aab}、\texttt{aabb}、\texttt{bb}。）
\begin{itemize}
\item \textbf{文法类型}：每条右部都是 $aB$ 或 $a$ 形（$S \to aS$、$S \to bA$、$A \to bA$ 等右线性，$\varepsilon$ 允许）\imp{} \textbf{3 型（正规文法）}。
\item \textbf{语言类型}：\textbf{正规语言}（正规式 $a^*b^*$，可画 DFA）。
\end{itemize}

\begin{qbox}
{\large\bfseries 第 6 题}\quad $G[S]$：$S \to aSb \mid bSa \mid SS \mid \varepsilon$。求 $L(G)$，并判断文法类型与语言类型；$G$ 是否二义文法？
\end{qbox}

\textbf{解}：$aSb$、$bSa$ 成对放一个 $a$、一个 $b$，$SS$ 拼接两段，$\varepsilon$ 是出口：
\[
L(G)=\{w \in \{a,b\}^* \mid w\text{ 中 }a\text{ 与 }b\text{ 的个数相等}\}
\]
\begin{itemize}
\item \textbf{正向}：对推导长度归纳——每条产生式都保持 $a$、$b$ 数目相等且 $a$、$b$ 可交错（如 \texttt{ab}、\texttt{ba}、\texttt{aabb}、\texttt{abab}、\texttt{abba}…）。
\item \textbf{反向}：任取 $a$、$b$ 等数的 $w$：若 $w$ 有真前缀也是等数的，$w = uv$，由 $S \Rightarrow SS \Rightarrow^* uv$；否则首字符与末字符必相异，$w = aw'b$ 或 $bw'a$ 且 $w'$ 等数，由 $S \Rightarrow aSb$（或 $bSa$）$\Rightarrow^* w$。归纳可得全部等数串，含 $\varepsilon$。
\item \textbf{文法类型}：$S \to SS$ 右部含两个非终结符 \imp{} 非 3 型 \imp{} \textbf{2 型}。
\item \textbf{语言类型}：\textbf{非正规}（等数要计数，正规语言泵引理可证不可行），是上下文无关语言。
\item \textbf{二义性}：\textbf{是}。句子 \texttt{abab} 有两棵语法树：$S \Rightarrow SS \Rightarrow (ab)(ab)$ 与 $S \Rightarrow aSb \Rightarrow a(ba)b$ \imp{} 二义文法。
\end{itemize}

\begin{qbox}
{\large\bfseries 第 7 题}\quad $G[S]$：$S \to 0S0 \mid 1S1 \mid 0 \mid 1$。求 $L(G)$，并判断文法类型与语言类型。
\end{qbox}

\textbf{解}：$0S0$、$1S1$ 在两头对称地加一对相同符号，$0$、$1$ 是奇数长度回文的中心：
\[
L(G)=\{w \in \{0,1\}^* \mid w = w^R \}
\]
（即 $0/1$ 上的\textbf{奇数长度回文串}：\texttt{0}、\texttt{1}、\texttt{010}、\texttt{111}、\texttt{00100}、\texttt{01010}…）
\begin{itemize}
\item \textbf{正向}：对推导长度归纳——对称包裹保持回文，长度每次 $+2$，出口长度 $1$ \imp{} 奇数长度回文。
\item \textbf{反向}：任一奇数长度回文 $w$，中心一个字符用 $0 \mid 1$，去掉中心后 $w = x x^R$，从外向内逐对用 $0S0/1S1$ 包裹即可。
\item \textbf{文法类型}：$0S0$、$1S1$ 右部非右线性 \imp{} \textbf{2 型}。
\item \textbf{语言类型}：\textbf{非正规}（回文要``记住前半段''，DFA 无此能力），是上下文无关语言。
\end{itemize}

\section{题型4：推导与语法树——短语、直接短语、句柄、素短语、最左素短语（第 8 题）}

\begin{qbox}
{\large\bfseries 第 8 题}\quad $G[E]$：$E \to E{+}T \mid E{-}T \mid T$，$T \to T{*}F \mid T/F \mid F$，$F \to (E) \mid i$。\\
对于句型 \textbf{E+T*F-i}：(1) 给出最右推导；(2) 画出语法树；(3) 写出所有短语、直接短语和句柄；(4) 写出所有素短语和最左素短语。
\end{qbox}

\textbf{解 (1) 最右推导}（每步都替换当前\textbf{最右}非终结符）：
\[
\begin{aligned}
E &\Rightarrow E-T        &&(E \to E-T)\\
  &\Rightarrow E-F        &&(T \to F)\\
  &\Rightarrow E-i        &&(F \to i)\\
  &\Rightarrow E+T-i      &&(E \to E+T)\\
  &\Rightarrow E+T*F-i    &&(T \to T*F)
\end{aligned}
\]

\textbf{解 (2) 语法树}：

\begin{center}
\begin{tikzpicture}[
  level distance=9.5mm,
  every node/.style={font=\small},
  level 1/.style={sibling distance=34mm},
  level 2/.style={sibling distance=13mm},
  level 3/.style={sibling distance=9mm},
  level 4/.style={sibling distance=12mm},
  edge from parent/.style={draw,thick}]
\node {$E$}
  child { node {$E$}
    child { node {$E$} }
    child { node {$+$} }
    child { node {$T$}
      child { node {$T$} }
      child { node {$*$} }
      child { node {$F$} }
    }
  }
  child { node {$-$} }
  child { node {$T$}
    child { node {$F$}
      child { node {$i$} }
    }
  };
\end{tikzpicture}
\end{center}

（根 $E \to E-T$；左 $E \to E+T$，其右 $T \to T*F$；根右 $T \to F \to i$。未展开的\textbf{叶型非终结符}：左下角的 $E$、$T$、$F$。）

\textbf{解 (3) 短语、直接短语、句柄}

\textbf{短语} = 语法树中任一子树的全部叶子从左到右组成的串：
\begin{itemize}
\item \textbf{E+T*F-i}（根 $E$ 的子树）
\item \textbf{E+T*F}（$E \to E+T$ 的子树）
\item \textbf{T*F}（$T \to T*F$ 的子树）
\item \textbf{i}（$T \to F \to i$ 的子树）
\end{itemize}

\textbf{直接短语} = 高度为 1 的子树（产生式右部\textbf{原样}出现在句型中，一步推导所得）：
\begin{itemize}
\item \textbf{T*F}（$T \to T*F$，$T$、$*$、$F$ 三符号都直接出现在句型中）
\item \textbf{i}（$F \to i$）
\end{itemize}

\textbf{句柄} = 最左直接短语 = \textbf{T*F}（也恰好是 (1) 中最右推导\textbf{最后一步}所用产生式 $T \to T*F$ 的右部，两法互验）。

\textbf{解 (4) 素短语、最左素短语}

\textbf{素短语} = 本身是短语、\textbf{至少含一个终结符}、且内部\textbf{不含更小短语}的串。把 (3) 的四个短语逐个检查：
\begin{itemize}
\item \textbf{i}：含终结符，内部无更小短语 \imp{} \textbf{素短语} \ok
\item \textbf{T*F}：含终结符 $*$，内部只有叶型非终结符 $T$、$F$（都不是短语）\imp{} \textbf{素短语} \ok
\item \textbf{E+T*F}：内部含更小短语 \textbf{T*F} \imp{} 不是素短语 \bad
\item \textbf{E+T*F-i}：内部含更小短语 \imp{} 不是素短语 \bad
\end{itemize}
\imp{} \textbf{素短语：T*F、i}；\textbf{最左素短语 = T*F}（最左边出现的素短语）。

\begin{tipbox}[句柄 vs 最左素短语]
本句型中句柄 = 最左素短语 = \textbf{T*F}，属巧合；二者一般不同——\textbf{句柄}是最左直接短语（可以是单个非终结符，不含终结符也可，如专题习题中的 $H \to G$ 的 \texttt{H}），\textbf{素短语}必须\textbf{至少含一个终结符}。规范归约每步归约\textbf{句柄}；\textbf{算符优先分析}每步归约\textbf{最左素短语}（见《期末复习专题》题型1 易错点 5）。
\end{tipbox}

\begin{warnbox}
\begin{enumerate}
\item \textbf{``E+T'' 不是短语}：它虽然是产生式 $E \to E+T$ 的右部、也原样出现在句型开头，但该子树的两个子结点 $E$、$T$ 都\textbf{继续展开了}，这棵子树实际推出的串是``E+T*F''。子树的``叶子串''才是短语——``E+T''只是一步产生式的形状，不是任何子树的叶子串。
\item \textbf{未展开的叶型非终结符（句型中的 $E$、$T$、$F$）不是短语}，它们只是子树叶子串的组成部分（对比《期末复习专题》题型4 习题2.10 修订警告）。
\item 句柄只有一个：\textbf{最左}直接短语；规范归约第一步就是把 \textbf{T*F} 归约为 $T$。
\end{enumerate}
\end{warnbox}

\section{\texorpdfstring{题型5：消除左递归（第 9--11 题：1 直接 + 2 间接）}{题型5：消除左递归（第 9–11 题：1 直接 + 2 间接）}}

\begin{qbox}
{\large\bfseries 第 9 题（直接左递归）}\quad $G[E]$：$E \to E{+}T \mid E{-}T \mid T$，$T \to T{*}F \mid T/F \mid F$，$F \to (E) \mid i$。消除文法的左递归。
\end{qbox}

\textbf{解}：$E$、$T$ 各含直接左递归（$F$ 无递归，保持不变）。按公式 $A \to A\alpha_1 \mid A\alpha_2 \mid \beta \Rightarrow A \to \beta A',\ A' \to \alpha_1 A' \mid \alpha_2 A' \mid \varepsilon$：
\begin{itemize}
\item $E \to E{+}T \mid E{-}T \mid T$：$\alpha_1=\text{“+T”}$、$\alpha_2=\text{“-T”}$、$\beta=\text{“T”}$ \imp{} $E \to TE'$，$E' \to {+}TE' \mid {-}TE' \mid \varepsilon$
\item $T \to T{*}F \mid T/F \mid F$：$\alpha_1=\text{“*F”}$、$\alpha_2=\text{“/F”}$、$\beta=\text{“F”}$ \imp{} $T \to FT'$，$T' \to {*}FT' \mid /FT' \mid \varepsilon$
\end{itemize}
结果：
\[
\begin{aligned}
E  &\to T E'\\
E' &\to + T E' \mid - T E' \mid \varepsilon\\
T  &\to F T'\\
T' &\to * F T' \mid / F T' \mid \varepsilon\\
F  &\to ( E ) \mid i
\end{aligned}
\]

\textbf{验证等价}：原 $E$ 可推得 $T_1 \pm T_2 \pm \cdots \pm T_n$（$n \ge 1$ 个 $T$ 由 $+$、$-$ 连接）；新文法 $E \to TE'$ 与 $E' \to \pm TE'$ 恰好生成同样的``1 个 $T$ + 若干个 $\pm T$''结构，$E' \Rightarrow \varepsilon$ 对应``只有一个 $T$''。$T$ 同理（$F$ 后接若干 ${*}F$、$/F$）。原语言无空串，新文法中 $E'$、$T'$ 的 $\varepsilon$ 仅供停止递归、不会产生额外终结串 \imp{} \textbf{语言等价}。

\begin{tipbox}[多条 $\alpha$、$\beta$ 要一一对应]
若 $A \to A\alpha_1 \mid A\alpha_2 \mid \beta_1 \mid \beta_2$，则 $A \to \beta_1 A' \mid \beta_2 A'$，$A' \to \alpha_1 A' \mid \alpha_2 A' \mid \varepsilon$。例如 $G[S]$：$S \to Sa \mid Sbc \mid b \mid c$ \imp{} \textbf{$S \to bS' \mid cS'$，$S' \to aS' \mid bcS' \mid \varepsilon$}（$\beta$ 有几条就写几条，别漏 $\varepsilon$，见《期末复习专题》题型5 易错点）。
\end{tipbox}

\begin{qbox}
{\large\bfseries 第 10 题（间接左递归·两非终结元）}\quad $G[S]$：$S \to Aa \mid b$，$A \to Sb \mid c$。消除左递归。
\end{qbox}

\textbf{解}：$S \Rightarrow Aa \Rightarrow Sba$，$S$（与 $A$）\textbf{间接左递归}。按非终结符顺序 $(S,\ A)$ 依次``把\textbf{前面}的产生式代入后面者，再消直接左递归''：
\begin{enumerate}
\item $S$ 的右部不以 $S$ 开头（无直接左递归），且无前序非终结符可代入 \imp{} $S \to Aa \mid b$ 不变。
\item $A$ 的右部含前序非终结符 $S$：代入 $S \to Aa \mid b$ 得 $A \to (Aa \mid b)b \mid c$，即 \textbf{$A \to Aab \mid bb \mid c$}（已化为直接左递归）。
\item 消直接左递归：$\alpha=\text{“ab”}$、$\beta_1=\text{“bb”}$、$\beta_2=\text{“c”}$ \imp{} $A \to bbA' \mid cA'$，$A' \to abA' \mid \varepsilon$。
\end{enumerate}
结果：
\[
\begin{aligned}
S  &\to A a \mid b\\
A  &\to b b A' \mid c A'\\
A' &\to a b A' \mid \varepsilon
\end{aligned}
\]

\textbf{验证等价}：代入消元——原 $A = Sb \mid c = (Aa \mid b)b \mid c = Aab \mid bb \mid c$ \imp{} $A=(bb \mid c)(ab)^*$，$S=Aa \mid b$；新文法给出的 $A$、$S$ 正是同一方程。核对小串：\texttt{b} \ok（$S \to b$）、\texttt{ca} \ok（$A \Rightarrow c \Rightarrow S \Rightarrow Aa \Rightarrow ca$）、\texttt{bba} \ok（$A \Rightarrow bb \Rightarrow bba$）、\texttt{bbaba} \ok（$A \Rightarrow bbab \Rightarrow bbaba$），且改写后无任何左递归（$A$ 的右部以 $b$、$c$ 打头，不再回到 $S$）\imp{} \textbf{语言等价}。

\begin{qbox}
{\large\bfseries 第 11 题（间接左递归·三非终结元成环）}\quad $G[S]$：$S \to Qc \mid c$，$Q \to Rb \mid b$，$R \to Sa \mid a$。消除左递归。
\end{qbox}

\textbf{解}：存在间接左递归环 $S \Rightarrow Qc \Rightarrow Rbc \Rightarrow Sabc$（$S$、$Q$、$R$ 同在一个圈上）。按非终结符顺序 $(S,\ Q,\ R)$ 执行``代入 + 消直接左递归''：
\begin{enumerate}
\item \textbf{$S$}：无前序；右部以 $Q$、$c$ 开头，无直接左递归 \imp{} $S \to Qc \mid c$ 不变。
\item \textbf{$Q$}：前序 $S$ 不出现在 $Q$ 的右部中 \imp{} 无事可做；无直接左递归 \imp{} $Q \to Rb \mid b$ 不变。
\item \textbf{$R$}：右部含前序 $S$，先代入 $S \to Qc \mid c$：$R \to Qca \mid ca \mid a$；再代入前序 $Q \to Rb \mid b$：\textbf{$R \to Rbca \mid bca \mid ca \mid a$}（已化为直接左递归）。
\item 消直接左递归：$\alpha=\text{“bca”}$、$\beta_1=\text{“bca”}$、$\beta_2=\text{“ca”}$、$\beta_3=\text{“a”}$ \imp{} $R \to bcaR' \mid caR' \mid aR'$，$R' \to bcaR' \mid \varepsilon$。
\end{enumerate}
结果：
\[
\begin{aligned}
S  &\to Q c \mid c\\
Q  &\to R b \mid b\\
R  &\to b c a R' \mid c a R' \mid a R'\\
R' &\to b c a R' \mid \varepsilon
\end{aligned}
\]

\textbf{验证等价}：代入消元——$R = Sa \mid a$ 代入 $Q = Rb \mid b$ 得 $Q = Sab \mid ab \mid b$，再代入 $S = Qc \mid c$ 得 $S = Sabc \mid abc \mid bc \mid c$ \imp{} $L(S)=(abc \mid bc \mid c)(abc)^*$。新文法：$R = (bca \mid ca \mid a)(bca)^*$，$Q = Rb \mid b$，$S = Qc \mid c$。核对小串：\texttt{c} \ok、\texttt{bc} \ok、\texttt{abc}（$a \cdot bc$）\ok、\texttt{cabc}（$ca \cdot bc$）\ok、\texttt{bcabc}（$bca \cdot bc$）\ok、\texttt{abcabc}（$a \cdot bca \cdot bc$）\ok 与 $(abc \mid bc \mid c)(abc)^*$ 逐一相同。改写后各右部均不以能推回自身的非终结符开头 \imp{} \textbf{无左递归、语言等价}。

\begin{tipbox}[间接左递归的通用套路（题型5 解题步骤）]
\begin{enumerate}
\item 定一个非终结符顺序（如 $S,\ Q,\ R$）；
\item 对第 $i$ 个非终结符：把 $j < i$ 的产生式\textbf{代入}其右部（直到右部只以 $i$ 或更后面的非终结符开头）；
\item 若此时出现直接左递归，套公式消去；
\item 顺序不同（如 $Q,\ R,\ S$）结果形状不同但\textbf{都等价}，写出``代入 + 消直接左递归''两个动作即可。
\item 注：原文法 $S \to Qc \mid c$ 等是\textbf{左线性 3 型}文法——左线性文法同样可能左递归、不能直接做自顶向下分析，消左递归正是为 LL 分析做准备（见《期末复习专题》题型5 知识点）。
\end{enumerate}
\end{tipbox}

\section{模拟卷·答案速查}

\begin{center}
\renewcommand{\arraystretch}{1.3}
\setlength{\tabcolsep}{6pt}
\begin{tabular}{@{}c p{13.6cm}@{}}
\toprule
\textbf{题号} & \textbf{答案要点} \\
\midrule
1 & 拓广文法 $Z \to \#E\#$；\FVname{}：$Z=\{\#\}$、$F=\{\neg,(,a\}$、$T=\{\wedge,\neg,(,a\}$、$E=\{\vee,\wedge,\neg,(,a\}$；\LVname{}：$Z=\{\#\}$、$F=\{\neg,),a\}$、$T=\{\wedge,\neg,),a\}$、$E=\{\vee,\wedge,\neg,),a\}$；表含 $\#$ 行列，仅 $( \doteq )$、$\# \doteq \#$ 两处同级、每格至多一种关系、无冲突 \imp{} \textbf{是算符优先文法} \\
2 & $S \to aA \mid \varepsilon$，$A \to aS$；\textbf{3 型}（写 $S \to aaS \mid \varepsilon$ 只是 2 型） \\
3 & $S \to AB$，$A \to aAb \mid ab$，$B \to cB \mid \varepsilon$；\textbf{2 型}（语言非正规） \\
4 & $S \to aSb \mid aS \mid a$；\textbf{2 型}（语言非正规） \\
5 & $L = \{a^i b^j \mid i,j \ge 0\} = a^* b^*$；文法 \textbf{3 型}、语言\textbf{正规} \\
6 & $L = \{w \in \{a,b\}^* \mid \#a = \#b\}$；文法 \textbf{2 型}、语言\textbf{非正规}、\textbf{二义} \\
7 & $L=\{0,1\}$ 上奇数长度回文；文法 \textbf{2 型}、语言\textbf{非正规} \\
8 & 短语：\textbf{E+T*F-i}、\textbf{E+T*F}、\textbf{T*F}、\textbf{i}；直接短语：\textbf{T*F}、\textbf{i}；\textbf{句柄 = T*F}；素短语：\textbf{T*F}、\textbf{i}，\textbf{最左素短语 = T*F} \\
9 & $E \to TE'$，$E' \to {+}TE' \mid {-}TE' \mid \varepsilon$；$T \to FT'$，$T' \to {*}FT' \mid /FT' \mid \varepsilon$；$F$ 不变 \\
10 & $S \to Aa \mid b$；$A \to bbA' \mid cA'$，$A' \to abA' \mid \varepsilon$ \\
11 & $S \to Qc \mid c$；$Q \to Rb \mid b$；$R \to bcaR' \mid caR' \mid aR'$，$R' \to bcaR' \mid \varepsilon$ \\
\bottomrule
\end{tabular}
\end{center}

\newpage
\begin{center}
{\LARGE\bfseries 期末复习强化卷 — 文法$\rightleftarrows$语言 8 题 + LL(1) 1 题（附答案与解析）}
\end{center}

\begin{infobox}
\textbf{用途}：期末冲刺第二份自测卷，共 9 题——\textbf{文法 $\to$ 语言 4 道、语言 $\to$ 文法 4 道、LL(1) 1 道}。全部为新题：参考《第二章PPT练习-完整题解》的题型与难度出题、\textbf{不重复其原题}；每题附\textbf{答案}与\textbf{解析}。\\
\textbf{配套}：《期末复习-模拟卷》·《期末复习专题》·《第二章PPT练习-完整题解》·《02-文法和语言》·《04-语法分析-自顶向下》\\
\textbf{做题方法}：文法$\rightleftarrows$语言题按``读递归模式 / 定骨架 \imp{} 写产生式 \imp{} 正反向验证 \imp{} 判类型''走（语言描述与文法写法均不唯一，多写法/多解见各题）；LL(1) 题按``FIRST \imp{} FOLLOW \imp{} SELECT \imp{} 判定 \imp{} 表''完整链条作答。
\end{infobox}

\section{\texorpdfstring{文法 $\to$ 语言，并判断类型（第 1--4 题）}{文法→语言，并判断类型（第 1–4 题）}}

\begin{qbox}
{\large\bfseries 第 1 题}\quad $G[S]$：$S \to Sa \mid Sb \mid c$。求 $L(G)$，并判断文法类型与语言类型。
\end{qbox}

\textbf{答案}（多种等价写法）：

\begin{center}
\renewcommand{\arraystretch}{1.25}
\begin{tabular}{c p{11.2cm}}
\toprule
\textbf{形式} & \textbf{写法} \\
\midrule
正规式 & $L(G)=c(a \mid b)^*$ \\
集合式 & $\{c\,w \mid w \in \{a,b\}^*\}$ \\
文字 & c 后随任意个 a、b 的串（含单独的 \texttt{c}） \\
例子 & \ok：\texttt{c}、\texttt{ca}、\texttt{cb}、\texttt{cab}、\texttt{caa}；\bad：\texttt{a}、\texttt{acb} \\
\bottomrule
\end{tabular}
\end{center}

文法 \textbf{3 型（左线性）}、语言\textbf{正规}。

\textbf{解析}：
\begin{itemize}
\item $S \to c$ 是出口；$S \to Sa$、$S \to Sb$ 每步在已有串\textbf{右端}追加一个字符（递归非终结符在右部\textbf{最左端}）。
\item \textbf{正向}：每用一次 $Sa/Sb$ 多一个 $a/b$ \imp{} 生成 $c$ 后接任意 $a/b$ 串；\textbf{反向}：任意 $c w$，把 $w$ 的字符从右往左逐个用 $Sa/Sb$ 追回，最后 $S \to c$ 收口 \imp{} 可推出。
\item \textbf{边界}：$w = \varepsilon$ \imp{} \texttt{c} \ok。
\item \textbf{类型}：三条产生式都是 $A \to Ba \mid a$ 形（\texttt{Sa} = $B \cdot a$，\texttt{Sb} = $B \cdot b$，\texttt{c} = $a$）\imp{} \textbf{3 型（左线性正规文法）}；语言有正规式 $c(a \mid b)^*$ \imp{} 正规。
\end{itemize}

\begin{warnbox}[方向别判反]
$S \to Sa$ 的递归非终结符在\textbf{左}、终结符在\textbf{右}——这是\textbf{左线性}（$A \to Ba \mid a$），不是右线性。左线性同属 3 型，但卷面写``右线性''就错了（对照专题 2.1 知识点对 3 型的两种形式）。
\end{warnbox}

\begin{qbox}
{\large\bfseries 第 2 题}\quad $G[S]$：$S \to aSbc \mid abc$。求 $L(G)$，并判断文法类型与语言类型。
\end{qbox}

\textbf{答案}（多种等价写法）：

\begin{center}
\renewcommand{\arraystretch}{1.25}
\begin{tabular}{c p{11.2cm}}
\toprule
\textbf{形式} & \textbf{写法} \\
\midrule
集合式 & $L(G)=\{a^n (bc)^n \mid n \ge 1\}$ \\
文字 & a 与 bc 块 1:1 配对、a 全在前（$n \ge 1$） \\
例子 & \ok：\texttt{abc}、\texttt{aabcbc}、\texttt{aaabcbcbc}；\bad：\texttt{a}、\texttt{ab}、\texttt{aabbc} \\
\bottomrule
\end{tabular}
\end{center}

文法 \textbf{2 型}、语言\textbf{非正规}。

\textbf{解析}：
\begin{itemize}
\item 出口 \texttt{abc}（$n=1$）；每套一层 $aSbc$ 在最前加一个 $a$、最后加一对 $bc$ \imp{} $a$ 与 $(bc)$ 块 $1{:}1$ 配对。
\item \textbf{正向}：归纳即得全部 $a^n(bc)^n$；\textbf{反向}：串的首字符必是 $a$、末两字符必是 \texttt{bc}，剥掉一层得 $a^{n-1}(bc)^{n-1}$，归纳落到 \texttt{abc}。
\item \textbf{边界}：$n \ge 1$ \imp{} $\varepsilon$、\texttt{a}、\texttt{ab} 都推不出 \ok。核对 $n=2$：$S \to aSbc \to a(abc)bc = \texttt{aabcbc}$ \ok。
\item \textbf{类型}：$S \to aSbc$ 右部非右线性，左部为单个非终结符 \imp{} \textbf{2 型}；$a$ 与 $bc$ 数目需配对计数 \imp{} 语言\textbf{非正规}。
\end{itemize}

\begin{tipbox}[\texorpdfstring{和 $a^n b^n c^n$ 的区别（关键）}{和 a^n b^n c^n 的区别（关键）}]
本语言的后两段是\textbf{交织}的 $(bc)^n$，一条 $S \to aSbc$ 就能配平 \imp{} \textbf{上下文无关}；而 $a^n b^n c^n$ 两段要分别配平、CFG 做不到，必须 1 型文法才能描述（对照题解第 4 题 (b)）。串形别看走眼：\texttt{aabcbc} $\in L$，\texttt{aabbc} $\notin L$。
\end{tipbox}

\begin{qbox}
{\large\bfseries 第 3 题}\quad $G[S]$：$S \to S;S \mid a$。求 $L(G)$，并判断文法类型与语言类型；$G$ 是否二义文法？
\end{qbox}

\textbf{答案}（多种等价写法）：

\begin{center}
\renewcommand{\arraystretch}{1.25}
\begin{tabular}{c p{11.2cm}}
\toprule
\textbf{形式} & \textbf{写法} \\
\midrule
正规式 & $L(G)=a( ; a)^*$ \\
文字 & 由分号分隔的 a 组成的非空串 \\
例子 & \ok：\texttt{a}、\texttt{a;a}、\texttt{a;a;a}；\bad：$\varepsilon$、\texttt{;}、\texttt{a;;a} \\
\bottomrule
\end{tabular}
\end{center}

文法 \textbf{2 型（非 3 型）}、语言\textbf{正规}、\textbf{$G$ 是二义文法}。

\textbf{解析}：
\begin{itemize}
\item $S \to a$ 出口；$S \to S;S$ 把两段用 \texttt{;} 拼接 \imp{} 任意 $\ge 1$ 个 $a$ 用 \texttt{;} 串起。
\item \textbf{正向}：逐次用 $S;S$ 拼接即得 $a( ; a)^*$；\textbf{反向}：对 $a;a;\cdots;a$，任取一个分号用 $S \to S;S$ 二分，递归到全是 $a$ \imp{} 只产生这类串。
\item \textbf{类型}：$S \to S;S$ 右部既不是 $aB$ 形也不是 $Ba$ 形（$S$、\texttt{;}、$S$ 三符号）\imp{} 非 3 型，左部单个非终结符 \imp{} \textbf{2 型}；语言 $a( ; a)^*$ 可用正规式表达 \imp{} \textbf{正规}——``语言正规 $\ne$ 文法 3 型''的又一实例。
\item \textbf{二义}：句子 \texttt{a;a;a} 有两棵语法树——左分 \texttt{((S;S);S)} 与右分 \texttt{(S;(S;S))} \imp{} \textbf{二义文法}（对分号连接链的结合方式没有规定）。
\end{itemize}

\begin{qbox}
{\large\bfseries 第 4 题}\quad $G[S]$：$S \to aSc \mid B$，$B \to bB \mid \varepsilon$。求 $L(G)$，并判断文法类型与语言类型。
\end{qbox}

\textbf{答案}（多种等价写法）：

\begin{center}
\renewcommand{\arraystretch}{1.25}
\begin{tabular}{c p{11.2cm}}
\toprule
\textbf{形式} & \textbf{写法} \\
\midrule
集合式 & $L(G)=\{a^n b^m c^n \mid n \ge 0,\ m \ge 0\}$ \\
文字 & 外层 a、c 等数配对，中间夹任意个 b（可为 0） \\
例子 & \ok：$\varepsilon$、\texttt{b}、\texttt{ac}、\texttt{abc}、\texttt{abbc}、\texttt{aacc}；\bad：\texttt{ca}、\texttt{abb} \\
\bottomrule
\end{tabular}
\end{center}

文法 \textbf{2 型}、语言\textbf{非正规}。

\textbf{解析}：
\begin{itemize}
\item $S \to aSc$ 每层配一对 $a$、$c$；出口 $S \to B$ 转入 $B$ 段，$B \to bB \mid \varepsilon$ 产出任意个 $b$（含 0 个）。串形：外层 $a \cdots c$ 夹着中间任意 $b$ 串。
\item \textbf{正向}：用 $k$ 次 $aSc$ 再 $S \to B$、$B$ 产 $m$ 个 $b$ \imp{} $a^k b^m c^k$；\textbf{反向}：串若以 $a$ 开头必以 $c$ 结尾（$n \ge 1$），剥一层递归；若无 $a$ 则整串是 $B$ 段（$b^m$ 或 $\varepsilon$）\imp{} 归纳成立。
\item \textbf{边界}：\textbf{$\varepsilon \in L(G)$}（$S \to B \to \varepsilon$）；$m$ 可为 $0$（\texttt{ac}、\texttt{aacc}…）；$n$ 可为 $0$（纯 $b$ 串 \texttt{b}、\texttt{bb}、$\varepsilon$）。
\item \textbf{类型}：$S \to aSc$ 非右线性 \imp{} \textbf{2 型}；$a$、$c$ 计数配对 \imp{} 语言\textbf{非正规}（中间 $b$ 串不限长只是``填充''，不影响判定）。
\end{itemize}

\begin{warnbox}[出口不能图省事]
若写 $S \to aSc \mid b \mid \varepsilon$：内层每次只能贡献\textbf{一个} $b$，得到的是``两边 $a$、$c$ 等数、中间至多一个 $b$''——\texttt{abc} 能推，但 \texttt{abbc} 推不出 \bad。$b^m$ 要\textbf{任意多个}，必须独立非终结符 $B$ 来递归。
\end{warnbox}

\section{\texorpdfstring{语言 $\to$ 文法，并判断类型（第 5--8 题，答案均不唯一）}{语言→文法，并判断类型（第 5–8 题，答案均不唯一）}}

\begin{qbox}
{\large\bfseries 第 5 题}\quad 写一文法 $G$，使 $L(G)=\{w \in \{a,b\}^* \mid w \text{ 以 ab 结尾}\}$，并判断文法类型。
\end{qbox}

\textbf{答案}（不唯一）：
\[
\begin{aligned}
S &\to aS \mid bS \mid aB\\
B &\to b
\end{aligned}
\]
\textbf{另一种}：$S \to aS \mid bS \mid ab$（更简，但右部 \texttt{ab} 是两个终结符 \imp{} 只是 \textbf{2 型}）。

\textbf{解析}：
\begin{itemize}
\item 拆两段：\textbf{前缀任意}（$S \to aS \mid bS$ 在前面逐字符加）+ \textbf{固定尾 \texttt{ab}}（$S \to aB \to ab$ 落底）。
\item \textbf{正向}：$S \to aS/bS$ 反复加前缀，最后 $S \to aB$ 产出结尾的 $ab$ \imp{} $(a \mid b)^* ab$；\textbf{反向}：$w = x \cdot ab$，$x$ 为空直接 $S \to aB$；$x = c\, x'$ 则 $S \to cS$ 递归剥首字符，剥完接 $aB$。
\item \textbf{边界}：最小串 \texttt{ab} \ok；$\varepsilon$、\texttt{a}、\texttt{b}、\texttt{ba} 均推不出 \ok（$S$ 每条路径都必须经过 $aB$）。
\item \textbf{类型}：全部 $A \to aB \mid a$ 形 \imp{} \textbf{3 型（正规文法）}；语言正规（3 状态 DFA：记住是否刚读过 $a$）。
\end{itemize}

\begin{qbox}
{\large\bfseries 第 6 题}\quad 写一文法 $G$，使 $L(G)=\{a^n b a^n \mid n \ge 0\}$（中间一个 $b$、两边 $a$ 等数：\texttt{b}、\texttt{aba}、\texttt{aabaa}…），并判断文法类型。
\end{qbox}

\textbf{答案}（不唯一）：
\[ S \to aSa \mid b \]
\textbf{另一种}：$S \to A$，$A \to aAa \mid b$（多包一层，语言相同）。

\textbf{解析}：
\begin{itemize}
\item $b$ 是中心，$aSa$ 每层在\textbf{两头同时}各加一个 $a$ \imp{} 两边永远等数。
\item \textbf{正向}：$n=0$ 用 $S \to b$；$n$ 层包裹 \imp{} $a^n b a^n$；\textbf{反向}：串首 $a$ 尾 $a$、中间恰一个 $b$ \imp{} 剥掉外层一对 $a$ 递归，最终落到 $b$。
\item \textbf{边界}：$n=0$ \imp{} \texttt{b} \ok；\texttt{ab}、\texttt{ba}、\texttt{aab} 推不出 \ok。
\item \textbf{核对}：$n=2$：$S \Rightarrow aSa \Rightarrow a(aSa)a \Rightarrow a(aba)a = \texttt{aabaa}$ \ok。
\item \textbf{类型}：$aSa$ 非右线性，左部单个非终结符 \imp{} \textbf{2 型}；两边 $a$ 数必须相等 \imp{} 语言\textbf{非正规}（DFA 记不住左边有几个 $a$）。
\end{itemize}

\begin{qbox}
{\large\bfseries 第 7 题}\quad 写一文法 $G$，使 $L(G)=\{a^m b^n c^{m+n} \mid m \ge 0,\ n \ge 0\}$（$c$ 的个数 = $a$、$b$ 个数之\textbf{和}），并判断文法类型。
\end{qbox}

\textbf{答案}（不唯一）：
\[
\begin{aligned}
S &\to aSc \mid T\\
T &\to bTc \mid \varepsilon
\end{aligned}
\]
\textbf{另一种}：$S \to aSc \mid bTc \mid \varepsilon$，$T \to bTc \mid \varepsilon$（把``转入 T 段''并成一条，语言相同）。

\textbf{解析}：
\begin{itemize}
\item 读结构：外层由 $S \to aSc$ 每层配一对 $a$、$c$（共 $m$ 层）；落到 $T$ 后内层由 $T \to bTc$ 每层配一对 $b$、$c$（共 $n$ 层）；$T \to \varepsilon$ 允许 $m$、$n$ 取 $0$。
\item \textbf{正向}：$m$ 次 $aSc$ 后 $S \to T$，$n$ 次 $bTc$ 后 $T \to \varepsilon$ \imp{} $a^m b^n c^{n} c^{m} = a^m b^n c^{m+n}$ \ok；\textbf{反向}：串首若为 $a$ 必尾为 $c$，剥一层外层递归；转入 $T$ 段后同理剥 $b \cdots c$。
\item \textbf{核对}：$m=1,\ n=1$ \imp{} $S \Rightarrow aSc \Rightarrow aTc \Rightarrow abTcc$ \imp{} \texttt{abcc}（$c^2 = 1+1$）\ok；$m=0,\ n=1$ \imp{} \texttt{bc} \ok；$m=1,\ n=0$ \imp{} \texttt{ac} \ok；$m=n=0$ \imp{} \textbf{$\varepsilon$} \ok（$S \to T \to \varepsilon$，\textbf{$\varepsilon$ 在语言中}）。
\item \textbf{类型}：$S \to aSc$ 非右线性 \imp{} \textbf{2 型}；$m+n$ 两段计数相加配对 \imp{} 语言\textbf{非正规}。
\end{itemize}

\begin{qbox}
{\large\bfseries 第 8 题}\quad 写一文法 $G$，使 $L(G)=\{w \in \{0,1\}^* \mid |w| \text{ 是 3 的倍数}\}$（含 $\varepsilon$），并判断文法类型。
\end{qbox}

\textbf{答案}（不唯一）：
\[
\begin{aligned}
S &\to 0A \mid 1A \mid \varepsilon\\
A &\to 0B \mid 1B\\
B &\to 0S \mid 1S
\end{aligned}
\]
\textbf{另一种}：语言等价写成 $((0 \mid 1)^3)^*$；文法的状态命名/转移顺序互换（如 A、B 对调并同步改转移）都等价。

\textbf{解析}：
\begin{itemize}
\item 按``已读长度 mod 3''设三个状态：\textbf{S}（余 0，也是接受状态）、\textbf{A}（余 1）、\textbf{B}（余 2）。每读一个字符状态 $+1$ 并循环：$S \to 0A/1A$，$A \to 0B/1B$，$B \to 0S/1S$；余 0 时可用 $S \to \varepsilon$ 结束。
\item \textbf{正向}：每步消耗一个字符、状态 $S \to A \to B \to S$ 循环，能结束（回到 $S$ 用 $\varepsilon$）当且仅当长度是 3 的倍数；\textbf{反向}：把串按每 3 个字符分段，每段三步从 $S$ 绕回 $S$，最后 $S \to \varepsilon$ 收口。
\item \textbf{核对}：$\varepsilon$（$S \to \varepsilon$）\ok、\texttt{011}（$S \Rightarrow 0A \Rightarrow 01B \Rightarrow 011S \Rightarrow 011$）\ok、\texttt{111000} \ok；\texttt{01}、\texttt{1111} 推不出（$A$、$B$ 没有 $\varepsilon$ 或直接收尾的产生式，走不回 $S$ 就断不了）\ok。
\item \textbf{类型}：每条右部都是 $aB$ / $a$（或 $\varepsilon$）形 \imp{} \textbf{3 型（正规文法）}；语言正规。
\end{itemize}

\section{LL(1) 文法：FIRST / FOLLOW / SELECT / 判定 / 预测分析表（第 9 题）}

\begin{qbox}
{\large\bfseries 第 9 题}\quad 已知文法 $G[E]$：
\[ \begin{aligned}
E  &\to T E'\\
E' &\to + T E' \mid \varepsilon\\
T  &\to F T'\\
T' &\to * F T' \mid \varepsilon\\
F  &\to ( E ) \mid i
\end{aligned} \]
(1) 求各非终结符的 \FRname 集；(2) 求各非终结符的 \FOname 集；(3) 求各产生式的 SELECT 集；(4) 判断 $G$ 是否为 LL(1) 文法，并说明理由；(5) 构造 LL(1) 预测分析表。
\end{qbox}

\textbf{解 (1) \FRname 集}（先标记可空：$F$、$T$ 不可空；$T'$、$E'$ 可空）：

\begin{center}
\renewcommand{\arraystretch}{1.25}
\begin{tabular}{c l p{8.6cm}}
\toprule
\textbf{非终结符} & \textbf{\FRname} & \textbf{依据} \\
\midrule
$F$  & $\{(,\ i\}$            & $F \to (E)$、$F \to i$ 直接加入 \\
$T'$ & $\{*,\ \varepsilon\}$  & $T' \to *FT'$ 加 $*$；$T' \to \varepsilon$ 使 $T'$ 可空 \\
$T$  & $\{(,\ i\}$            & $T \to FT'$：\FR{F} 不含 $\varepsilon$，直接取 \FR{F} \\
$E'$ & $\{+,\ \varepsilon\}$  & $E' \to +TE'$ 加 $+$；$E' \to \varepsilon$ 使 $E'$ 可空 \\
$E$  & $\{(,\ i\}$            & $E \to TE'$：\FR{T} 不含 $\varepsilon$，直接取 \FR{T} \\
\bottomrule
\end{tabular}
\end{center}

\textbf{解 (2) \FOname 集}（从开始符号 $E$ 出发，沿``某 $X$ 后面跟什么''传播；\textbf{可空则再并上后面的 \FOname}——本题两条传播链是易错点）：

\begin{center}
\renewcommand{\arraystretch}{1.25}
\begin{tabular}{c l p{8.2cm}}
\toprule
\textbf{非终结符} & \textbf{\FOname} & \textbf{依据} \\
\midrule
$E$  & $\{\#,\ )\}$            & （1）$E$ 是开始符号 \imp{} $\#$；（2）$F \to (E)$ 中 $E$ 后是 $)$ \imp{} $)$ \\
$E'$ & $\{\#,\ )\}$            & $E \to TE'$ 中 $E'$ 在右部末端 \imp{} $\FO{E'} = \FO{E}$ \\
$T$  & $\{+,\#,\ )\}$          & $E \to TE'$：$T$ 后跟 $E'$ \imp{} 取 $\FR{E'} \setminus \{\varepsilon\} = \{+\}$；又 \textbf{$E'$ 可空} \imp{} 再并 $\FO{E} = \{\#, )\}$ \\
$T'$ & $\{+,\#,\ )\}$          & $T \to FT'$ 中 $T'$ 在末端 \imp{} $\FO{T'} = \FO{T}$ \\
$F$  & $\{*,+,\#,\ )\}$        & $T \to FT'$：$F$ 后是 $T'$ \imp{} 取 $\FR{T'} \setminus \{\varepsilon\} = \{*\}$；又 \textbf{$T'$ 可空} \imp{} 再并 $\FO{T} = \{+,\#, )\}$ \\
\bottomrule
\end{tabular}
\end{center}

\textbf{解 (3) SELECT 集}（非 $\varepsilon$ 产生式取 $\mathrm{FIRST}(\text{右部})$；$\varepsilon$ 产生式取 $\mathrm{FOLLOW}(\text{左部})$）：

\begin{center}
\renewcommand{\arraystretch}{1.25}
\begin{tabular}{l l l}
\toprule
\textbf{产生式} & \textbf{计算} & \textbf{SELECT} \\
\midrule
$E \to T E'$     & \FR{TE'}     & $\{(,\ i\}$ \\
$E' \to + T E'$  & \FRname        & $\{+\}$ \\
$E' \to \varepsilon$ & \FO{E'}  & $\{\#,\ )\}$ \\
$T \to F T'$     & \FR{FT'}     & $\{(,\ i\}$ \\
$T' \to * F T'$  & \FRname        & $\{*\}$ \\
$T' \to \varepsilon$ & \FO{T'}  & $\{+,\#,\ )\}$ \\
$F \to ( E )$    & \FRname        & $\{(\}$ \\
$F \to i$        & \FRname        & $\{i\}$ \\
\bottomrule
\end{tabular}
\end{center}

\textbf{解 (4) 判定}：相同左部的产生式两两求 SELECT 集交集：
\begin{itemize}
\item $E$：只有 $E \to TE'$ 一条，无交集可算；
\item $E'$：$SELECT(E' \to +TE') \cap SELECT(E' \to \varepsilon) = \{+\} \cap \{\#, )\} = \emptyset$；
\item $T$：只有 $T \to FT'$ 一条，无交集可算；
\item $T'$：$SELECT(T' \to {*}FT') \cap SELECT(T' \to \varepsilon) = \{*\} \cap \{+,\#, )\} = \emptyset$；
\item $F$：$SELECT(F \to (E)) \cap SELECT(F \to i) = \{( \} \cap \{i\} = \emptyset$。
\end{itemize}
全部为空 \imp{} \textbf{相同左部的任意两产生式 SELECT 集交集为空，所以 $G$ 是 LL(1) 文法}。

\textbf{解 (5) LL(1) 预测分析表}（$M[\text{非终结符},\ \text{终结符}]$，每格恰好一个产生式，空格为报错）：

\begin{center}
\renewcommand{\arraystretch}{1.3}
\begin{tabular}{c|cccccc}
\toprule
        & $($          & $i$       & $+$            & $*$             & $)$           & $\#$ \\
\midrule
\textbf{E}  & $E \to TE'$ & $E \to TE'$ &                &                 &               & \\
\textbf{E'} &             &             & $E' \to +TE'$  &                 & $E' \to \varepsilon$ & $E' \to \varepsilon$ \\
\textbf{T}  & $T \to FT'$ & $T \to FT'$ &                &                 &               & \\
\textbf{T'} &             &             & $T' \to \varepsilon$ & $T' \to *FT'$ & $T' \to \varepsilon$ & $T' \to \varepsilon$ \\
\textbf{F}  & $F \to (E)$ & $F \to i$   &                &                 &               & \\
\bottomrule
\end{tabular}
\end{center}

（逐格来源：$E$、$T$ 按 \FRname 填；$E'$、$T'$ 的 $\varepsilon$ 行按 \FOname 填；$F$ 按两条 \FRname 分列。）

\begin{warnbox}
\begin{enumerate}
\item \textbf{\FO{T}、\FO{F} 的``可空传播''}是全题最大失分点：$E'$ 可空才让 \texttt{+、\#、)} 进入 \FO{T}；$T'$ 可空才让 \texttt{+、\#、)} 进入 \FO{F}。漏传播 \imp{} $\mathrm{SELECT}(T' \to \varepsilon)$ 算错 \imp{} 判定跟着错。
\item \textbf{$\varepsilon$ 产生式的 SELECT 用 \FOname 不是 \FRname}（$\FR{\varepsilon} = \emptyset$ 没用）。
\item 判定写法：同左部产生式的 SELECT 交集\textbf{逐对列出}、全部为空，再下结论。
\end{enumerate}
\end{warnbox}

\section{强化卷·答案速查}

\begin{center}
\renewcommand{\arraystretch}{1.3}
\setlength{\tabcolsep}{6pt}
\begin{tabular}{@{}c p{13.6cm}@{}}
\toprule
\textbf{题号} & \textbf{答案要点} \\
\midrule
1 & $L = c(a \mid b)^*$（多写法见题）；\textbf{3 型（左线性）}、正规 \\
2 & $L = \{a^n (bc)^n \mid n \ge 1\}$（多写法见题）；\textbf{2 型}、非正规（区别于 $a^n b^n c^n$） \\
3 & $L = \{a( ; a)^*\}$（多写法见题）；\textbf{2 型}（非 3 型）、语言正规、\textbf{二义} \\
4 & $L = \{a^n b^m c^n \mid n,m \ge 0\}$（\textbf{含 $\varepsilon$}，多写法见题）；\textbf{2 型}、非正规（出口须经独立 $B$） \\
5 & $S \to aS \mid bS \mid aB$，$B \to b$（另解 $S \to aS \mid bS \mid ab$，2 型）；\textbf{3 型}、正规 \\
6 & $S \to aSa \mid b$（另解 $S \to A$，$A \to aAa \mid b$）；\textbf{2 型}、非正规 \\
7 & $S \to aSc \mid T$，$T \to bTc \mid \varepsilon$（另解 $S \to aSc \mid bTc \mid \varepsilon$）；\textbf{2 型}、非正规（\textbf{含 $\varepsilon$}） \\
8 & $S \to 0A \mid 1A \mid \varepsilon$，$A \to 0B \mid 1B$，$B \to 0S \mid 1S$（另解 $((0 \mid 1)^3)^*$）；\textbf{3 型}、正规（mod 3 状态机） \\
9 & FIRST/FOLLOW/SELECT 见解(1)(2)(3)；同左部 SELECT 交集全空 \imp{} \textbf{是 LL(1) 文法}（判定见解(4)）；表见解(5) \\
\bottomrule
\end{tabular}
\end{center}

\end{document}
